\scanpage{71}
Calculer $F$ à l'aide de sa restriction à $\mathbb{K}$ comme
\[
F(X)\simeq\unsure{\mathbb R|R}^0F(L)\simeq\operatorname{Ker}(F(I)\longrightarrow F(J))
\]
Ceci dit, \deleted{\unsure{si}} on a
\diagram{\begin{tikzcd}[ampersand replacement=\&,column sep=large]
F(X)\otimes_B F(Y)\arrow[r,"{?}"]\arrow[d] \& F(X*Y)\arrow[d]\\
F(I)\otimes_B F(J)\arrow[r,"{\substack{\text{connu}\\\text{par } F\vert\mathbb{K}\ \text{et sa}\\\text{structure mult.}}}"'] \& F(I*J)
\end{tikzcd}}

Or on voit que \deleted{si} $*$ transforme des couples de $\Sigma$-monos en un $\Sigma$-mono (résulte immédiatement des formules de Dold Puppe) donc la $2^{\text{ème}}$ flèche verticale est un mono d'où la détermination unique de ?. (Le fait qu'il soit indépendant des choix faits est immédiat, et aussi qu'il satisfait aux compatibilités d'ass. et de comm. si sa restriction à $\mathbb{K}$ y satisfait. Il reste donc à examiner l'axiome d'unité, \deleted{qui est que} il faut donc se donner
\[
\deleted{F(\Psi_0)}\mathbb 1_{B'}\longrightarrow F(\Psi_0]\xrightarrow{\sim}Z^0(C^{\bullet})
\]
et on voit que l'axiome signifie que c'est une \uline{unité} pour la structure de ``$k$-$\otimes$-Alg'' de $F(\deleted{\unsure{\Phi}}\Psi^0)$ (qui est une ``$k$-$\otimes$-Alg ass. comm'' \unsure{de tte}).
% margin note written sideways
\margin{\uline{NB} Donc on voit qu'on n'a pas besoin de l'existence de $\operatorname{Ker}(d^0:C^0\longrightarrow C^1)$ il suffit de dire qu'on donne $\mathbb 1_B\longrightarrow C^0$ dont le composé avec $C^0\longrightarrow C^1$ est nul, et qui joue le rôle d'une unité \deleted{\unsure{de $F(C^0$}} en un sens qu'on \unsure{explicitera|conjecture}}
% archive sheet number: 70
